\chapter{Definitions}

The forms described in this chapter are \emph{Scheme procedures} called
at the meta-level --- in a \texttt{.scm} source file or at the Scheme
REPL --- not expressions in the VNB object language.  Each call installs
axioms or macetes into the current theory as a side effect; nothing is
proved.  The distinction between meta-level definition and object-level
proof bears repeating: \texttt{declare-structure}, \texttt{def-constant},
and their companions manipulate the theory data structure directly, while
proofs are constructed separately using the proof commands of
Chapter~\ref{ch:proofs}.

%% =========================================================
\section{Constants and predicates}
\label{sec:def-constant}

A \emph{defined constant} is a new constant symbol introduced by a
collection of characterizing axioms.  It differs from a plain
\texttt{theory-add-axiom!} call only in bookkeeping: the axioms are
recorded under a separate \texttt{definitions} slot in the theory record,
so the system can distinguish what was \emph{postulated} from what was
\emph{introduced by definition}.

\begin{VNBcode}
(def-constant 'NAME
  '(axiom-name-1 formula-1)
  '(axiom-name-2 formula-2)
  ...)
\end{VNBcode}

\begin{sloppypar}Each extra argument is a two-element list \texttt{(axiom-name formula)}.
The formulas are installed in the global theorem and macete tables exactly
as \texttt{theory-add-axiom!} would install them, so they are immediately
available to proof commands and macetes.\end{sloppypar}

\textbf{Existence and uniqueness are the user's responsibility.}  The
system accepts the definition without checking that the described object
exists or is unique.  For constants where this matters (e.g., a least
element), you should prove existence and uniqueness before invoking
\texttt{def-constant}, or at least note the obligation explicitly.

Example---introducing the first uncountable ordinal $\omega_1$:
\begin{VNBcode}
(def-constant 'OMEGA1
  (list 'omega1-is-ordinal
        (parse-string "OMEGA1 in ORD"))
  (list 'omega1-is-uncountable
        (parse-string "not(countable(OMEGA1))"))
  (list 'omega1-is-least-unc
        (parse-string "forall([alpha in ORD],
          not(countable(alpha)) implies OMEGA1 <= alpha)")))
\end{VNBcode}

To display all defined constants and their characterizing axioms:
\begin{VNBcode}
(display-definitions)
\end{VNBcode}

%% =========================================================
\section{Predicates}
\label{sec:def-predicate}

An $n$-ary \emph{defined predicate} is introduced by \texttt{def-predicate}.
It installs a single characterising axiom --- an \texttt{iff} --- that
unfolds the predicate in terms of its arguments.

\begin{VNBcode}
(def-predicate 'PRED-NAME '(p1 p2 ...)  body)
\end{VNBcode}

Installs one axiom named \texttt{PRED-NAME}:
\begin{VNBsnip}
forall([p1, p2, ...], PRED-NAME(p1, p2, ...) iff body)
\end{VNBsnip}

The axiom is immediately available via \texttt{(ta 'PRED-NAME)} in a proof.

\paragraph{Example: Cauchy sequence in a metric space.}
\begin{VNBcode}
(def-predicate 'IS-CAUCHY-SEQ '(X f)
  '(AND (IS-METRIC-SPACE X)
        (IN f (FUN NN (CARRIER X)))
        (FORALL eps
          (IMPLIES (AND (IN eps RR) (< 0 eps))
            (FORSOME N (IN N NN)
              (FORALL m (IMPLIES (IN m NN)
                (FORALL n (IMPLIES (IN n NN)
                  (IMPLIES (AND (<= N m) (<= N n))
                    (< ((D X) (f m) (f n)) eps)))))))))))
\end{VNBcode}

This installs:
\begin{VNBsnip}
forall([X, f],
  IS-CAUCHY-SEQ(X, f) iff
    IS-METRIC-SPACE(X)
    and f in fun(NN, CARRIER(X))
    and forall([eps in RR], 0 < eps implies
          forsome([N in NN],
            forall([m in NN, n in NN], N <= m and N <= n implies
              D(X)(f(m), f(n)) < eps)))
\end{VNBsnip}

%% =========================================================
\section{Recursive definitions}
\label{sec:recursive-defs}

Two Scheme helpers introduce function constants by recursion and install
the defining equations as axioms.  Both rely on the same
\texttt{theory-add-definition!} mechanism as \texttt{def-constant}
(Section~\ref{sec:def-constant}), so every axiom is immediately
available in proofs.

\begin{prop}[Primitive recursion on $\mathit{NN}$]\label{prop:nn-rec}
Let $A \in \mathit{SET}$ and let
\[
  \mathcal{F}(A) \;=\; \bigcup_{n \in \mathit{NN}}\,\mathit{FUN}(\mathit{ORD\text{-}SEGMENT}(n),\, A)
\]
be the set of all functions defined on a finite initial segment of $\mathit{NN}$
with values in $A$ (including the empty function at $n = 0$, since
$\mathit{ORD\text{-}SEGMENT}(0) = \emptyset$).
For any $\mathcal{E} \in \mathit{FUN}(\mathcal{F}(A),\, A)$ there is a unique
$f \in \mathit{FUN}(\mathit{NN},\, A)$ such that
\[
  f(n) \;=\; \mathcal{E}\!\bigl(f \!\restriction\! \mathit{ORD\text{-}SEGMENT}(n)\bigr)
  \qquad \text{for every } n \in \mathit{NN}.
\]
\end{prop}

\begin{remark}
This is the special case of Proposition~\ref{prop:tfrec-set} with
$U = \mathit{NN}$; it follows from the well-ordering of $\mathit{NN}$ by the
standard least-counterexample argument.
Since $\mathit{ORD\text{-}SEGMENT}(n)$ is the integer interval $\{0, 1, \ldots, n-1\}$,
the extension function $\mathcal{E}$ sees the entire history
$f(0), f(1), \ldots, f(n-1)$ when computing $f(n)$.
The \emph{predecessor form} --- $f(0) = b$ and
$f(\mathrm{succ}(n)) = s(n,\, f(n))$ for $n \in \mathit{NN}$ --- is the special case
where $\mathcal{E}$ only inspects the last value; this covers sums, products,
and most concrete recursive definitions.
\end{remark}

\subsection{\texttt{def-by-nn-recursion}: primitive recursion on \texttt{NN}}

Defines a function $F(p_1, \ldots, p_k, n)$ by two-clause recursion on the
last argument $n \in \mathit{NN}$, with arbitrary extra parameters
$p_1, \ldots, p_k$:
\[
  F(p_1,\ldots,p_k,0) = b, \qquad
  F(p_1,\ldots,p_k,\mathrm{succ}(n)) = s(p_1,\ldots,p_k,n,F(p_1,\ldots,p_k,n)).
\]

\begin{VNBcode}
(def-by-nn-recursion 'F
  '(p1 p2 ...)      ; extra parameters ('() for a plain NN->A function)
  base-val          ; F(p1,p2,..., 0) = base-val
  '(n val)          ; step variables
  succ-expr)        ; F(p1,p2,..., succ(n)) = succ-expr
                    ;   (val stands for F(p1,p2,...,n))
\end{VNBcode}

Two axioms are installed:
\begin{VNBsnip}
F-zero: forall([p1,p2,...], F(p1,p2,...,0) = base-val)
F-succ: forall([p1,p2,..., n in NN],
          F(p1,p2,...,succ(n)) = succ-expr[val:=F(p1,p2,...,n)])
\end{VNBsnip}

\paragraph{Example: factorial (no parameters).}
\begin{VNBcode}
(def-by-nn-recursion 'FACT '()
  1                           ; FACT(0) = 1
  '(n val) '(* (succ n) val)) ; FACT(succ(n)) = succ(n) * FACT(n)
\end{VNBcode}

\paragraph{Example: power of 2 (no parameters).}
\begin{VNBcode}
(def-by-nn-recursion 'POW2 '()
  1                       ; POW2(0) = 1
  '(n val) '(* 2 val))   ; POW2(succ(n)) = 2 * POW2(n)
\end{VNBcode}

\paragraph{Example: partial sum (one parameter).}
\begin{sloppypar}For $f \in \mathit{FUN}(\mathit{NN}, A)$, define
$\mathrm{SUM}(f, n) = \sum_{k=0}^{n} f(k)$:\end{sloppypar}
\begin{VNBcode}
(def-by-nn-recursion 'SUM '(f)
  '(f 0)                          ; SUM(f, 0) = f(0)
  '(n val) '(+ val (f (succ n)))) ; SUM(f, succ(n)) = SUM(f,n) + f(succ(n))

; Installs SUM-zero: forall([f], SUM(f,0) = f(0))
; Installs SUM-succ: forall([f, n in NN], SUM(f,succ(n)) = SUM(f,n) + f(succ(n)))
\end{VNBcode}

\subsection{\texttt{def-by-ord-recursion}: transfinite recursion on \texttt{ORD}}

For functions on all ordinals, use the three-clause schema
(Proposition~\ref{prop:tfrec-class}, Section~\ref{sec:ord-recursion}):
\begin{itemize}[label=$\circ$, leftmargin=0.9cm]
\item \textbf{Base} ($\xi = 0$):\quad $F(0) = b$
\item \textbf{Successor} ($\xi = \mathrm{succ\_ORD}(\alpha)$):\quad
      $F(\mathrm{succ\_ORD}(\alpha)) = s(\alpha,\, F(\alpha))$
\item \textbf{Limit} ($\mathrm{LIMIT\text{-}ORD}(\xi)$):\quad
      $F(\xi) = \ell\!\left(\xi,\, F\!\restriction\!\mathrm{ORD\text{-}SEGMENT}(\xi)\right)$
\end{itemize}

\begin{VNBcode}
(def-by-ord-recursion 'F
  base-val            ; F(0) = base-val
  '(alpha val)        ; successor: F(succ_ORD(alpha)) = succ-expr
  succ-expr           ;   where val stands for F(alpha)
  '(lambda)           ; limit: F(lambda) = lim-expr
  lim-expr)           ;   where lambda is the limit ordinal
\end{VNBcode}

Three axioms are installed:
\begin{VNBsnip}
F-zero:  F(0) = base-val
F-succ:  forall([alpha in ORD],
           F(succ_ORD(alpha)) = succ-expr[val:=F(alpha)])
F-limit: forall([lambda], LIMIT-ORD(lambda) implies F(lambda) = lim-expr)
\end{VNBsnip}

\paragraph{Example: identity function on ordinals.}
\begin{VNBcode}
(def-by-ord-recursion 'COPY-ORD
  0                                         ; COPY-ORD(0) = 0
  '(alpha val) '(succ_ORD val)              ; COPY-ORD(succ(a)) = succ(COPY-ORD(a))
  '(lam) '(SUP-ORD (ORD-SEGMENT lam)))      ; COPY-ORD(lam) = sup of predecessors
\end{VNBcode}

\paragraph{Example: ordinal addition $\alpha + \beta$ (fixing $\alpha$).}

Defining $A_\alpha(\beta) = \alpha + \beta$ as a function of $\beta$:
\begin{VNBcode}
(def-by-ord-recursion 'ORD-ADD
  'alpha-fixed                              ; ORD-ADD(0) = alpha-fixed
  '(beta val) '(succ_ORD val)              ; ORD-ADD(succ(b)) = succ(ORD-ADD(b))
  '(lam) '(SUP-ORD (ORD-SEGMENT lam)))     ; ORD-ADD(lam) = sup of ORD-ADD on lam
\end{VNBcode}

(Here \texttt{alpha-fixed} is the constant playing the role of $\alpha$.
A full definition would use nested ordinal recursion.)

%% =========================================================
\section{Mathematical structures}
\label{sec:declare-structure}

\subsection{The \texttt{declare-structure} form}

\texttt{declare-structure} declares a named mathematical structure at the
Scheme level.  It registers the structure, installs NTH-based accessor
macetes for every carrier set and every operation, and posts a
definitional axiom for the associated \texttt{IS-\textit{NAME}} predicate
in the current theory.  No quoting is needed.

\begin{VNBcode}
(declare-structure NAME
  (carriers C1 C2 ...)
  (op OPNAME (D1 D2 ...) RANGE)
  (constant CNAME SET))
\end{VNBcode}

\begin{itemize}[label = $\circ$, leftmargin=1cm]

\item \texttt{(carriers C1 C2 ...)}
  Names the carrier sets.  A structure with $m$ carriers has its $i$-th
  carrier at \texttt{nth($i$, s)} in any instance~\texttt{s}.

\item \texttt{(op NAME (D1 D2 ...) RANGE)}
  Declares a function-valued slot.  If more than one domain component is
  listed, the domain is \texttt{CARTESIAN(D1, D2, ...)};
  a single-element list gives a unary domain.  All carrier and op names
  in domain and range positions are automatically expanded to accessor
  applications.

\item \texttt{(constant NAME SET)}
  Declares a constant-valued slot: $\mathit{NAME}(s) \in \mathit{SET}(s)$.

\end{itemize}

Accessor indices are assigned sequentially: carriers occupy indices
$1, \ldots, m$ in declaration order; operations and constants follow
at indices $m+1, m+2, \ldots$.

\begin{remark}
A \texttt{declare-structure} declaration involves three distinct things,
which it is worth keeping clearly separated.
\begin{enumerate}
\item \textbf{The name} --- a symbol such as \texttt{RING} or
  \texttt{METRIC-SPACE}, used as a label in declarations and prose.
\item \textbf{An instance} --- a specific VNB list, which is an
  \emph{element of \textsc{set}} (a concrete set-theoretic object).
  For example,
  \[
    \mathtt{CC\text{-}MS} \;=\;
      [\,\mathtt{CC},\;\lambda([\mathtt{x},\mathtt{y}],\,
      \mathtt{magnitude}(\mathtt{x}-\mathtt{y}))\,]
  \]
  is a 2-element VNB list.  Internally it is \texttt{(LIST CC (VNB-LAMBDA
  (LIST x y) (magnitude (- x y))))}.  The accessors \texttt{X(s)} and
  \texttt{D(s)} are \texttt{nth(1,~s)} and \texttt{nth(2,~s)}.
\item \textbf{The associated class} --- the collection
  $\{\,s : \mathit{IS\text{-}NAME}(s)\,\}$ of \emph{all} VNB lists
  satisfying the structure's typing conditions.  \emph{This class
  shares the structure's name}: \texttt{declare-structure NAME}
  installs both the predicate \texttt{IS-NAME} and the class
  \texttt{NAME}, related by the axiom \texttt{NAME-class}:
  \[
    \forall s.\; s \in \mathit{NAME} \;\Leftrightarrow\; \mathit{IS\text{-}NAME}(s).
  \]
  The class is a proper class (defined by a predicate), not an
  element of \textsc{set}.  Naming the class is what makes the
  natural quantification idiom
  $\texttt{forall}([s \in \mathit{NAME}], \ldots)$
  work --- without it, one would have to write the unbounded form
  $\texttt{forall}([s], \mathit{IS\text{-}NAME}(s) \Rightarrow \ldots)$.
\end{enumerate}
This matches the standard textbook presentation in which a ring, say,
is defined as a 6-tuple $(R,+,\times,-,0,1)$ satisfying certain axioms:
the tuple is the instance (point~2) and the collection of all such tuples
is the class (point~3), now an addressable name.
\end{remark}

\subsection{What \texttt{declare-structure} installs}

\paragraph{Accessor macetes.}
For each carrier $C_i$ (at NTH index $i$) and each operation
$\mathit{op}_j$ (at NTH index $m+j$), a macete is installed that
rewrites the accessor term to the corresponding NTH projection:
\[
  C_i(s) \;\longmapsto\; \texttt{nth}(i,\; s)
  \qquad
  \mathit{op}_j(s) \;\longmapsto\; \texttt{nth}(m{+}j,\; s)
\]
The schema variable matches any expression, so the macete fires anywhere
in a goal.

\paragraph{IS-\textit{NAME} axiom.}
The predicate \texttt{IS-\textit{NAME}} is axiomatized as a conjunction
whose first conjunct asserts that $s$ is a VNB list of exactly the right
length, followed by typing conditions on each accessor:
\begin{align*}
  \forall s.\;\mathit{IS\text{-}NAME}(s) \;\Leftrightarrow\; &
  \mathtt{length}(s) = n \\
  &\land\; C_1(s) \in \mathit{SET} \land \cdots \land C_m(s) \in \mathit{SET}\\
  &\land\; \mathit{op}_1(s) \in \mathit{FUN}(\mathit{dom}_1, \mathit{rng}_1)
  \land \cdots
\end{align*}
where $n = m + k$ is the total number of components ($m$ carriers and $k$
operations/constants).  The length condition ties the predicate to the
instance side of the three-part picture: only a VNB list of exactly $n$
elements can satisfy \texttt{IS-\textit{NAME}}.

\begin{sloppypar}The axiom is named \texttt{IS-\textit{NAME}} and is immediately available
via \texttt{cmd-theorem-assumption} and \texttt{cmd-apply-macete}.\end{sloppypar}

\paragraph{\textit{NAME}-class axiom.}
A second axiom named \texttt{\textit{NAME}-class} ties the predicate to
the class:
\[
  \forall s.\; s \in \mathit{NAME} \;\Leftrightarrow\; \mathit{IS\text{-}NAME}(s).
\]
The symbol $\mathit{NAME}$ thereby denotes the proper class
$\{s \mid \mathit{IS\text{-}NAME}(s)\}$, ready to use as the
quantification range:
\begin{VNBsnip}
forall([r in RING], blah(r))             ; bounded over the ring class
forall([E in NORMEDSPACE], blah(E))      ; bounded over normed spaces
\end{VNBsnip}
which the parser expands to
$\forall s.\; s \in \mathit{NAME} \Rightarrow \ldots$; the
\texttt{\textit{NAME}-class} axiom then converts the antecedent to
the predicate form $\mathit{IS\text{-}NAME}(s)$ exactly when needed.

Characteristic axioms (associativity, distributivity, etc.) are
\emph{not} generated automatically; they are installed by separate
\texttt{theory-add-axiom!} calls after the declaration.

\subsection{Pre-loaded algebraic structures}
\label{sec:algebraic}

The file \texttt{algebraic.scm} is loaded automatically and declares the
following structures, each paired with its characteristic axioms.

\subsubsection{Semigroup}

\begin{VNBcode}
(declare-structure SEMIGROUP
  (carriers A)
  (op MUL (A A) A))
; A -> (NTH 1 s),  MUL -> (NTH 2 s)
\end{VNBcode}

One axiom:
\begin{VNBsnip}
semigroup-assoc:
  forall([s], IS-SEMIGROUP(s) implies
    forall([x in A(s), y in A(s), z in A(s)],
      MUL(s)(MUL(s)(x, y), z) = MUL(s)(x, MUL(s)(y, z))))
\end{VNBsnip}

\subsubsection{Monoid}

\begin{VNBcode}
(declare-structure MONOID
  (carriers A)
  (op MUL (A A) A)
  (constant E A))
; A -> (NTH 1 s),  MUL -> (NTH 2 s),  E -> (NTH 3 s)
\end{VNBcode}

Three axioms: \texttt{monoid-assoc}, \texttt{monoid-left-id}, \texttt{monoid-right-id}.
\begin{VNBsnip}
monoid-left-id:
  forall([s], IS-MONOID(s) implies
    forall([a in A(s)], MUL(s)(E(s), a) = a))

monoid-right-id:
  forall([s], IS-MONOID(s) implies
    forall([a in A(s)], MUL(s)(a, E(s)) = a))
\end{VNBsnip}

Two convenience lemmas are installed in \texttt{axioms.scm}:
\begin{description}
\item[\texttt{monoid-identity-in}]
  $\mathit{IS\text{-}MONOID}(m) \Rightarrow \mathit{E}(m) \in \mathit{A}(m)$.
\item[\texttt{monoid-carrier-closed-mul}]
  $\mathit{IS\text{-}MONOID}(m) \land a \in \mathit{A}(m) \land b \in \mathit{A}(m)
  \Rightarrow \mathit{MUL}(m)(a,b) \in \mathit{A}(m)$.
\end{description}

\subsubsection{Group}

\begin{VNBcode}
(declare-structure GROUP
  (carriers A)
  (op MUL (A A) A)
  (constant E A)
  (op INV (A) A))
; A -> (NTH 1 s),  MUL -> (NTH 2 s),  E -> (NTH 3 s),  INV -> (NTH 4 s)
\end{VNBcode}

\begin{sloppypar}Three axioms suffice (left-only): \texttt{group-assoc}, \texttt{group-left-id},
\texttt{group-left-inv}.\end{sloppypar}
\begin{VNBsnip}
group-left-inv:
  forall([s], IS-GROUP(s) implies
    forall([a in A(s)], MUL(s)(INV(s)(a), a) = E(s)))
\end{VNBsnip}

\subsubsection{Ring}

\begin{VNBcode}
(declare-structure RING
  (carriers A)
  (op ADD  (A A) A)
  (op MUL  (A A) A)
  (op NEG  (A)   A)
  (constant ZERO A)
  (constant ONE  A))
; A->1, ADD->2, MUL->3, NEG->4, ZERO->5, ONE->6
\end{VNBcode}

\begin{sloppypar}
Nine axioms: \texttt{ring-add-assoc}, \texttt{ring-add-comm},
\texttt{ring-add-left-id}, \texttt{ring-add-left-inv},
\texttt{ring-mul-assoc}, \texttt{ring-mul-left-id}, \texttt{ring-mul-right-id},
\texttt{ring-left-dist}, \texttt{ring-right-dist}.
\end{sloppypar}

$(A, \mathit{ADD}, \mathit{ZERO}, \mathit{NEG})$ is an abelian group;
$(A, \mathit{MUL}, \mathit{ONE})$ is a monoid; $\mathit{MUL}$ distributes
over $\mathit{ADD}$.

Two convenience lemmas are installed in \texttt{axioms.scm}:
\begin{description}
\item[\texttt{ring-zero-in}]
  $\mathit{IS\text{-}RING}(r) \Rightarrow \mathit{ZERO}(r) \in \mathit{A}(r)$.
\item[\texttt{ring-carrier-closed-add}]
  $\mathit{IS\text{-}RING}(r) \land a \in \mathit{A}(r) \land b \in \mathit{A}(r)
  \Rightarrow \mathit{ADD}(r)(a,b) \in \mathit{A}(r)$.
\end{description}

\subsubsection{Metric space}

\begin{VNBcode}
(declare-structure METRIC-SPACE
  (carriers X)
  (op D (X X) RR))
; X -> (NTH 1 s),  D -> (NTH 2 s)
\end{VNBcode}

Five axioms: \texttt{metric-pos}, \texttt{metric-self-zero},
\texttt{metric-zero-eq}, \texttt{metric-sym}, \texttt{metric-triangle}.
\begin{VNBsnip}
metric-triangle:
  forall([s], IS-METRIC-SPACE(s) implies
    forall([u in X(s), v in X(s), w in X(s)],
      D(s)(u, w) <= D(s)(u, v) + D(s)(v, w)))
\end{VNBsnip}

\subsection{Further structure definitions}

The same mechanism handles any arity of carriers and operations.  For
instance, a Banach space over $\mathbb{C}$:

\begin{VNBcode}
(declare-structure NORMEDSPACE
  (carriers VECTORS)
  (op SCALAR-TIMES (CC VECTORS) VECTORS)
  (op PLUS         (VECTORS VECTORS) VECTORS)
  (op MINUS        (VECTORS VECTORS) VECTORS)
  (op NORM         (VECTORS) RR))
; VECTORS->1, SCALAR-TIMES->2, PLUS->3, MINUS->4, NORM->5
\end{VNBcode}

This installs five accessor macetes and adds the following axiom:
\begin{VNBsnip}
forall([s],
  IS-NORMEDSPACE(s) iff
    VECTORS(s) in SET
    and SCALAR-TIMES(s) in fun(cartesian(CC, VECTORS(s)), VECTORS(s))
    and PLUS(s) in fun(cartesian(VECTORS(s), VECTORS(s)), VECTORS(s))
    and MINUS(s) in fun(cartesian(VECTORS(s), VECTORS(s)), VECTORS(s))
    and NORM(s) in fun(VECTORS(s), RR))
\end{VNBsnip}

A characteristic axiom relating \texttt{MINUS} to \texttt{PLUS} and
\texttt{SCALAR-TIMES} is installed separately.  Here the axiom is written
in S-expression form rather than string notation, because the string
parser reads \texttt{-1} as unary minus applied to \texttt{1} rather than
as a negative integer literal; the S-expression \texttt{-1} is a genuine
Scheme integer and avoids that ambiguity.

\begin{VNBcode}
(theory-add-axiom! *current-theory* 'normedspace-minus
  '(FORALL s
     (IMPLIES (IS-NORMEDSPACE s)
       (FORALL x (IMPLIES (IN x (VECTORS s))
         (FORALL y (IMPLIES (IN y (VECTORS s))
           (= ((MINUS s) x y)
              ((PLUS s) x ((SCALAR-TIMES s) -1 y))))))))))
\end{VNBcode}

A vector space over an arbitrary field:
\begin{VNBcode}
(declare-structure VECTORSPACE
  (carriers FIELD VECTORS)
  (op SCALAR-TIMES (FIELD VECTORS) VECTORS)
  (op PLUS         (VECTORS VECTORS) VECTORS)
  (constant ZERO VECTORS)
  (op NEG          (VECTORS) VECTORS))
; FIELD->1, VECTORS->2, SCALAR-TIMES->3, PLUS->4, ZERO->5, NEG->6
\end{VNBcode}

%%----------------------------------------------------------
\subsection{Structure specialization: \texttt{specialize-structure}}
\label{sec:specialize-structure}

Every characteristic axiom and every theorem proved about a structure class
has the generic shape
\[
  \forall s.\;\mathit{IS\text{-}NAME}(s) \;\Rightarrow\; P[s].
\]
Once a specific instance $h$ has been \emph{certified}---meaning a theorem
of the form $\mathit{IS\text{-}RING}(h)$ has been proved and installed---every
such generic statement can be specialized to $h$ by universal instantiation
followed by modus ponens, yielding the concrete statement $P[h]$.

\texttt{specialize-structure} automates this transport in bulk:

\begin{VNBcode}
(specialize-structure 'HARP 'RING 'harp-is-ring)
; => 17 theorems specialized for HARP.
\end{VNBcode}

The interactive short form is \texttt{(spec instance struct is-thm)}.

\paragraph{Arguments.}
\begin{itemize}[label=$\circ$, leftmargin=1cm]
\item \textit{instance-name} --- symbol naming the specific structure, e.g.\
  \texttt{HARP}.
\item \textit{struct-name} --- the structure class, e.g.\ \texttt{RING}.
  The predicate \texttt{IS-RING} is derived automatically.
\item \textit{is-thm-name} --- name of a theorem in the current theory
  whose formula is exactly $\mathit{IS\text{-}RING}(\mathit{HARP})$.
\end{itemize}

\paragraph{What gets specialized.}
Every entry in the global theorem table (axioms and proved theorems alike)
whose formula matches
\[
  \forall s.\;\mathit{IS\text{-}RING}(s) \;\Rightarrow\; P[s]
\]
--- with $s$ the bound variable appearing as the sole argument to
$\mathit{IS\text{-}RING}$ --- is specialized.  The result $P[\mathit{HARP}]$
is installed under the name
$\langle\textit{original\text{-}name}\rangle\texttt{-harp}$.
Axioms such as \texttt{ring-add-assoc} are included, so
\texttt{ring-add-assoc-harp} is immediately available after the call.

\paragraph{Soundness.}
No new axioms are introduced.  Each specialized theorem $P[h]$ follows
from the certified $\mathit{IS\text{-}RING}(h)$ and the generic
$\forall s.\;\mathit{IS\text{-}RING}(s)\Rightarrow P[s]$ by two standard
inferences: universal instantiation and modus ponens.
\texttt{specialize-structure} performs these in bulk without requiring a
separate interactive proof for each one.  The certification check at the
start of the call verifies that \textit{is-thm-name} has the expected form
before any specialization occurs; an incorrect theorem name raises an error.

\begin{remark}
Theorems about sub-structures are \emph{not} automatically transported.
For instance, theorems of the form
$\forall s.\;\mathit{IS\text{-}GROUP}(s)\Rightarrow\ldots$
are not touched by \texttt{(spec 'HARP 'RING ...)}.  To obtain them you
would need to certify $\mathit{IS\text{-}GROUP}(\mathit{HARP})$ separately
and call \texttt{spec} again for the \texttt{GROUP} class.
\end{remark}

\paragraph{Example: the complex numbers as a ring.}
The constant \texttt{CC-RING} is defined in \texttt{complex.scm} as the
6-tuple $[\mathit{CC},\,+,\,\times,\,-,\,0,\,1]$.  The axiom
\texttt{cc-is-ring} asserts $\mathit{IS\text{-}RING}(\mathit{CC\text{-}RING})$.

\begin{VNBcode}
(spec 'CC-RING 'RING 'cc-is-ring)
; => n theorems specialized for CC-RING.
; Installs: ring-add-assoc-cc-ring, ring-add-comm-cc-ring,
;           ring-mul-assoc-cc-ring, ring-left-dist-cc-ring, ...
\end{VNBcode}

A proof about $\mathbb{C}$ can then retrieve any of these directly:
\begin{VNBcode}
(ta 'ring-add-comm-cc-ring)
; adds  ADD(CC-RING)(x, y) = ADD(CC-RING)(y, x)  to the context
\end{VNBcode}

%% =========================================================
\section{\texttt{def-functoid}}
\label{sec:def-functoid}

\texttt{def-functoid} introduces a named functoid at the proof-checker
level.  It is a definitional device, not an object of the theory: the
name it installs can be applied and reasoned about, but there is no
formal quantifier ranging over functoids.

The general form is:
\begin{VNBcode}
(def-functoid name (x) body)
\end{VNBcode}
where \texttt{x} is a formal parameter ranging over classes and
\texttt{body} is a class-valued expression, typically built from
accessor macetes installed by \texttt{declare-structure}.

\paragraph{Example: forgetful functor NORMEDSPACE $\to$ METRIC-SPACE.}
A normed space carries an induced metric $d(x,y) = \|x - y\|$.
Since the domain is the proper class of all normed spaces, \texttt{fun}
cannot express this mapping and a functoid is required:

\begin{VNBcode}
(def-functoid NormedToMetric (S)
  [(VECTORS S),
   (lambda [[x, y] in cartesian((VECTORS S), (VECTORS S))],
     ((NORM S) ((MINUS S) x y)))])
\end{VNBcode}

The result is a tuple \texttt{[VECTORS(S), D]} matching the
\texttt{METRIC-SPACE} signature (carrier \texttt{X}, distance \texttt{D}).
That it satisfies \texttt{IS-METRIC-SPACE} is a theorem, not automatic:
the metric axioms follow from the norm axioms and the definition of
\texttt{MINUS}.

\paragraph{Example: forgetful functor NORMEDSPACE $\to$ VECTORSPACE.}
Dropping the norm recovers the underlying vector space:

\begin{VNBcode}
(def-functoid NormedToVector (S)
  [(VECTORS S), (SCALAR-TIMES S), (PLUS S), (MINUS S)])
\end{VNBcode}

\paragraph{Example: the carrier functor.}
\begin{VNBcode}
(def-functoid GroupCarrier (G) (CARRIER G))
\end{VNBcode}
This maps every group to its underlying carrier set.  The result is
always a set (\texttt{IS-GROUP} entails \texttt{CARRIER(G) in SET}),
even though the domain is a proper class.

\paragraph{Transfinite recursion and functoids.}
\begin{sloppypar}The transfinite recursion scheme
(Section~\ref{sec:ord-recursion}) produces a functoid, not a
set-theoretic function: its domain is the proper class $\mathit{ORD}$
of all ordinals, so the result cannot live in $\mathit{FUN}(\mathit{ORD}, B)$
for any set-valued $B$.  A typical pattern is to define the step
functor $\mathcal{E}$ as a functoid via \texttt{def-functoid} and
then invoke \texttt{def-by-ord-recursion}:\end{sloppypar}
The actual signature has six arguments: the new functoid name, the
base value, plus a binding spec and body for both the successor and
limit cases.  The successor body is evaluated with the bound name
referring to the recursive value at the predecessor; the limit body
is evaluated with the bound name referring to the partial functoid
restricted to all earlier ordinals.
\begin{VNBcode}
(def-by-ord-recursion 'MySeq
  'a0                                 ; base value at 0
  '(prev) '(succ_ORD prev)            ; successor: MySeq(succ alpha) = succ_ORD(prev)
  '(part) '(SUP-ORD (image part)))    ; limit:    MySeq(lambda) = sup of partial seq
; => installs axioms
;      MySeq(0) = a0
;      IN alpha ORD => MySeq(succ_ORD alpha) = succ_ORD(MySeq alpha)
;      LIMIT-ORD lambda => MySeq(lambda) = SUP-ORD(image-of MySeq below lambda)
\end{VNBcode}

%%% -----------------------------------------------------------------------
\section{Cardinality: \texttt{CARD}}
\label{sec:cardinality}

\texttt{CARD} is a functoid mapping every set $A$ to its cardinality
$\mathit{CARD}(A) \in \mathit{ORD}$.  Under the axiom of choice,
$\mathit{CARD}(A)$ is the least ordinal $\alpha$ that bijects with $A$
(the Hartogs--Scott construction).  We axiomatise the key properties
directly in \texttt{cardinality.scm}.

\begin{remark}
Because the domain of $\mathit{CARD}$ is the proper class $\mathit{SET}$,
it cannot be expressed as an element of $\mathit{FUN}(\mathit{SET}, \mathit{ORD})$
(that class is not a set).  It is therefore a functoid, on the same
footing as $\mathit{ORD\text{-}SEGMENT}$.
\end{remark}

\begin{remark}
$\mathit{CARD}$ is total over all classes, but characterised only when
its argument is a set.  The kernel does not define $\mathit{CARD}(A)$
for a proper class $A$: the term $\mathit{CARD}(\mathit{ORD})$, for
instance, is syntactically well-formed but no axiom constrains its
value.  This is consistent with the totality framing (every functoid
returns \emph{some} value), but means that any equational use of
$\mathit{CARD}(A)$ outside the established sethood-conditioned axioms
will not discharge.  Convention: only invoke $\mathit{CARD}$ on a
class you have proven to be a set.
\end{remark}

\subsection*{Axioms}

\begin{description}
\item[\texttt{card-in-ord}]
  $A \in \mathit{SET} \Rightarrow \mathit{CARD}(A) \in \mathit{ORD}$.

\item[\texttt{card-empty}]
  $\mathit{CARD}(\emptyset) = 0$.

\item[\texttt{card-insert}]
  $A \in \mathit{SET} \wedge x \notin A \;\Rightarrow\;
   \mathit{CARD}(A \cup \{x\}) = \mathit{succ_{ORD}}(\mathit{CARD}(A))$.\\
  Here $\{x\} = \mathit{PAIR}(x,x)$ (by the pairing axiom with $a = b = x$).

\item[\texttt{card-segment}]
  $n \in \mathit{NN} \;\Rightarrow\; \mathit{CARD}(\mathit{ORD\text{-}SEGMENT}(n)) = n$.\\
  Bridge between cardinality and the $\mathit{NN}$-indexed model of finite sets.

\item[\texttt{card-finite-bij}]
  $A \in \mathit{SET} \wedge \mathit{CARD}(A) \in \mathit{NN} \;\Rightarrow\;$
  there exists a bijection
  $\varphi : \mathit{ORD\text{-}SEGMENT}(\mathit{CARD}(A)) \to A$.\\
  (Existence follows from AC; stated as an axiom here.)

\item[\texttt{card-union-disjoint}]
  $A, B \in \mathit{SET},\;
   \mathit{CARD}(A), \mathit{CARD}(B) \in \mathit{NN},\;
   A \cap B = \emptyset \;\Rightarrow\;
   \mathit{CARD}(A \cup B) = \mathit{CARD}(A) + \mathit{CARD}(B)$.\\
  Finite additivity; the $+$ on the right is $\mathit{NN}$ addition.
\end{description}

\subsection*{Finiteness criterion}

A set $A$ is \emph{finite} iff $\mathit{CARD}(A) \in \mathit{NN}$.
For finite sets, \texttt{card-segment} says
$\mathit{CARD}(\mathit{ORD\text{-}SEGMENT}(n)) = n$, so the initial
segment $\{0, 1, \ldots, n-1\}$ has cardinality exactly $n$.  The
axiom \texttt{card-finite-bij} supplies an enumeration bijection for
any finite set, which is the key ingredient for defining unordered
finite sums in Section~\ref{sec:sequences}.

%%% -----------------------------------------------------------------------
\section{Monoid products and ring sums}
\label{sec:sequences}

\texttt{sequences.scm} defines two constants by primitive recursion on
$\mathit{NN}$:

\begin{itemize}
\item $\mathit{PROD\text{-}ORD}(m, f, n)$ --- product of $f(0) * f(1) * \cdots * f(n{-}1)$
  in a monoid $m$.
\item $\mathit{SUM}(r, f, n)$ --- sum $f(0) + f(1) + \cdots + f(n{-}1)$
  in a ring $r$.
\end{itemize}

Both are $\mathit{NN}$-indexed: the third argument $n$ counts how many
terms are included.  The empty product ($n = 0$) is the identity;
$\mathit{SUM}(r, f, 0) = \mathit{ZERO}(r)$.

\subsection{\texttt{PROD-ORD}: monoid product}

\paragraph{Signature.}
$\mathit{PROD\text{-}ORD}(m, f, n)$ where $m$ is a \textsc{Monoid},
$f \in \mathit{FUN}(\mathit{NN}, \mathit{A}(m))$, $n \in \mathit{NN}$.

\paragraph{Defining recursion.}
Installed by \texttt{def-by-nn-recursion} as two named axioms:

\begin{align*}
\mathtt{prod\text{-}ord\text{-}zero}: &\quad
  \forall m\, f.\;
  \mathit{PROD\text{-}ORD}(m, f, 0) = \mathit{E}(m) \\
\mathtt{prod\text{-}ord\text{-}succ}: &\quad
  \forall m\, f\, n \in \mathit{NN}.\;
  \mathit{PROD\text{-}ORD}(m, f, \mathit{succ}(n))
  = \mathit{PROD\text{-}ORD}(m, f, n) * f(n)
\end{align*}

where $*$ denotes $\mathit{MUL}(m)$.

\paragraph{Additional axioms.}
\begin{description}
\item[\texttt{prod-ord-type}]
  $\mathit{IS\text{-}MONOID}(m) \wedge f \in \mathit{FUN}(\mathit{NN}, \mathit{A}(m))
  \wedge n \in \mathit{NN} \;\Rightarrow\;
  \mathit{PROD\text{-}ORD}(m, f, n) \in \mathit{A}(m)$.
  Installed as an axiom for convenience; formally derivable from
  \texttt{prod-ord-zero}, \texttt{prod-ord-succ}, \texttt{monoid-identity-in},
  and \texttt{monoid-carrier-closed-mul} by $\mathit{NN}$-induction
  (see Section~\ref{sec:ex-ni-type}).

\item[\texttt{prod-ord-singleton}]
  $\mathit{IS\text{-}MONOID}(m) \wedge f \in \mathit{FUN}(\mathit{NN}, \mathit{A}(m))
  \;\Rightarrow\; \mathit{PROD\text{-}ORD}(m, f, 1) = f(0)$.\\
  (Derivable from \texttt{prod-ord-succ} at $n=0$, \texttt{prod-ord-zero},
  and \texttt{monoid-left-id}.)
\end{description}

\subsection{\texttt{SUM}: ring summation}

\texttt{SUM} is the special case of \texttt{PROD-ORD} for the additive
monoid of a ring---\textit{i.e.}, $\mathit{MUL} \to \mathit{ADD}$,
$\mathit{E} \to \mathit{ZERO}$---defined independently for direct use.

\paragraph{Signature.}
$\mathit{SUM}(r, f, n)$ where $r$ is a \textsc{Ring},
$f \in \mathit{FUN}(\mathit{NN}, \mathit{A}(r))$, $n \in \mathit{NN}$.

\paragraph{Defining recursion.}

\begin{align*}
\mathtt{sum\text{-}zero}: &\quad
  \forall r\, f.\; \mathit{SUM}(r, f, 0) = \mathit{ZERO}(r) \\
\mathtt{sum\text{-}succ}: &\quad
  \forall r\, f\, n \in \mathit{NN}.\;
  \mathit{SUM}(r, f, \mathit{succ}(n))
  = \mathit{SUM}(r, f, n) + f(n)
\end{align*}

\paragraph{Additional axioms.}
\begin{description}
\item[\texttt{sum-type}]
  $\mathit{IS\text{-}RING}(r) \wedge f \in \mathit{FUN}(\mathit{NN}, \mathit{A}(r))
  \wedge n \in \mathit{NN} \;\Rightarrow\; \mathit{SUM}(r, f, n) \in \mathit{A}(r)$.
  Installed as an axiom for convenience; formally derivable from
  \texttt{sum-zero}, \texttt{sum-succ}, \texttt{ring-zero-in},
  and \texttt{ring-carrier-closed-add} by $\mathit{NN}$-induction
  (Section~\ref{sec:ex-ni-type}).

\item[\texttt{sum-singleton}]
  $\mathit{IS\text{-}RING}(r) \wedge f \in \mathit{FUN}(\mathit{NN}, \mathit{A}(r))
  \;\Rightarrow\; \mathit{SUM}(r, f, 1) = f(0)$.

\item[\texttt{sum-left-scalar}]
  $a \in \mathit{A}(r) \wedge f \in \mathit{FUN}(\mathit{NN}, \mathit{A}(r))
  \wedge n \in \mathit{NN} \;\Rightarrow\;$\\
  $\mathit{SUM}(r,\, \lambda i.\, a \cdot f(i),\, n)
  = a \cdot \mathit{SUM}(r, f, n)$.\\
  (Derivable by $\mathit{NN}$-induction using \texttt{ring-left-dist}.)
\end{description}

\begin{remark}
The unordered finite sum $\sum_{x \in S} f(x)$ over an arbitrary finite
set $S$ (not indexed by $\mathit{NN}$) requires commutativity of
addition and a well-definedness proof (independence of the enumeration
bijection from \texttt{card-finite-bij}).  This more general
$\mathit{SUM\text{-}SET}$ operation is planned but not yet implemented.
\end{remark}

\begin{remark}
The connection $\mathit{SUM}(r, f, n) = \mathit{PROD\text{-}ORD}(m, f, n)$
where $m = [\mathit{A}(r),\, \mathit{ADD}(r),\, \mathit{ZERO}(r)]$ is the
additive monoid of $r$ is a theorem (by structural induction on $n$),
not an axiom.
\end{remark}

\paragraph{Typical proof pattern.}
To use \texttt{SUM} in a proof, unfold the recursion manually with
\texttt{(ta 'sum-succ)} and \texttt{(ta 'sum-zero)}, then close
numeric identities with \texttt{(rs)} or \texttt{(arith)}.

\begin{VNBcode}
;; Prove SUM(r, f, 2) = f(0) + f(1)
(sp (make-wff-from-string
  "forall([r, f], IS-RING(r) and f in FUN(NN, A(r)) implies
     SUM(r, f, 2) = ADD(r)(SUM(r, f, 1), f(1))"))
(ta 'sum-succ)    ; unfold outer step
(ta 'sum-singleton)
(rfl)
\end{VNBcode}

%% =========================================================
\section{Ring products}
\label{sec:ring-prod}

\texttt{algebraic.scm} defines the binary ring product, and
\texttt{sequences.scm} extends it to finite indexed products.
Both follow Path~2: the operations are \emph{total} functoids, defined
for any arguments; typing theorems fire conditionally when the inputs
are rings.

\subsection{Zero ring: \texttt{ZERO-RING}}

\texttt{ZERO-RING} is the terminal ring: carrier $\{0\}$, all operations
returning $0$, with $\mathit{ZERO} = \mathit{ONE} = 0$.

\begin{VNBsnip}
ZERO-RING = [{0},
             lambda([p, q], 0),   ; ADD
             lambda([p, q], 0),   ; MUL
             lambda([p], 0),      ; NEG
             0,                   ; ZERO
             0]                   ; ONE
\end{VNBsnip}

\begin{description}
\item[\texttt{zero-ring-def}]
  $\mathit{ZERO\text{-}RING} = [\{0\},\,
   \lambda(p,q).\,0,\; \lambda(p,q).\,0,\; \lambda(p).\,0,\; 0,\; 0]$.
\item[\texttt{zero-ring-is-ring}]
  $\mathit{IS\text{-}RING}(\mathit{ZERO\text{-}RING})$.
  (Taken as an axiom; the proof would verify the nine ring axioms for
  the trivial one-element ring.)
\end{description}

\begin{remark}
The zero ring is the terminal object in the category of rings: there is
exactly one ring homomorphism from \emph{any} ring $R$ to
$\mathit{ZERO\text{-}RING}$ (mapping everything to $0$).  It is also
the identity for \texttt{RING-PROD} up to isomorphism:
$\mathit{RING\text{-}PROD}(\mathit{ZERO\text{-}RING}, R) \cong R$.
\end{remark}

\subsection{Binary product: \texttt{RING-PROD}}

\texttt{RING-PROD(X,~Y)} is the product ring with carrier
$\mathit{CARTESIAN}(\mathit{A}(X),\,\mathit{A}(Y))$ and
componentwise operations.

\paragraph{Definition.}
Installed as a \texttt{def-functoid} rewrite:

\begin{VNBsnip}
RING-PROD(X, Y) =
  [CARTESIAN(A(X), A(Y)),
   lambda([p, q], [ADD(X)(nth(1,p), nth(1,q)),
                   ADD(Y)(nth(2,p), nth(2,q))]),
   lambda([p, q], [MUL(X)(nth(1,p), nth(1,q)),
                   MUL(Y)(nth(2,p), nth(2,q))]),
   lambda([p],    [NEG(X)(nth(1,p)),
                   NEG(Y)(nth(2,p))]),
   [ZERO(X), ZERO(Y)],
   [ONE(X),  ONE(Y)]]
\end{VNBsnip}

Elements of $\mathit{A}(\mathit{RING\text{-}PROD}(X,Y))$ are pairs $[a,b]$
with $a \in \mathit{A}(X)$ and $b \in \mathit{A}(Y)$;
all six operations act componentwise.

\paragraph{Accessor expansion.}
The \texttt{RING-PROD} macete chain is:
\[
  \mathit{A}(\mathit{RING\text{-}PROD}(X,Y))
  \;\longrightarrow\; \mathit{CARTESIAN}(\mathit{A}(X),\,\mathit{A}(Y))
\]
and similarly for \texttt{ADD}, \texttt{MUL}, \texttt{NEG}, \texttt{ZERO},
\texttt{ONE}: each expands by first applying the \texttt{RING-PROD}
rewrite and then the accessor macete.

\paragraph{Typing axiom.}
\begin{description}
\item[\texttt{ring-prod-is-ring}]
  $\mathit{IS\text{-}RING}(X) \wedge \mathit{IS\text{-}RING}(Y)
  \;\Rightarrow\; \mathit{IS\text{-}RING}(\mathit{RING\text{-}PROD}(X,Y))$.
  (Taken as an axiom.  The proof would verify that the componentwise
  operations satisfy each ring axiom, given that $X$ and $Y$ do.)
\end{description}

\begin{remark}[Totality]
\texttt{RING-PROD} is total: $\mathit{RING\text{-}PROD}(X,Y)$ is a
well-defined tuple for \emph{any} $X$ and $Y$.  When neither is a ring
--- when bongos are passed in --- the result is some tuple that satisfies
no ring axioms.  \texttt{ring-prod-is-ring} is silent on this case (its
antecedent is false), so no contradiction arises.  This is the VNB
tradition: totality and typing are separate concerns.
\end{remark}

\subsection{Finite indexed product: \texttt{RING-PROD-N}}

$\mathit{RING\text{-}PROD\text{-}N}(f, n)$ is the $n$-fold product ring
$f(0) \times f(1) \times \cdots \times f(n-1)$, defined by primitive
recursion on~$n$.

\paragraph{Signature.}
$f$ is any function on $\mathit{NN}$ (not required to map to rings;
the typing axiom states when the result is a ring).
$n \in \mathit{NN}$.

\paragraph{Defining recursion.}
Installed by \texttt{def-by-nn-recursion}:

\begin{align*}
\mathtt{ring\text{-}prod\text{-}n\text{-}zero}: &\quad
  \forall f.\; \mathit{RING\text{-}PROD\text{-}N}(f, 0) = \mathit{ZERO\text{-}RING} \\
\mathtt{ring\text{-}prod\text{-}n\text{-}succ}: &\quad
  \forall f,\, n \in \mathit{NN}.\;
  \mathit{RING\text{-}PROD\text{-}N}(f, \mathit{succ}(n))
  = \mathit{RING\text{-}PROD}(\mathit{RING\text{-}PROD\text{-}N}(f, n),\; f(n))
\end{align*}

\paragraph{Typing axiom.}
\begin{description}
\item[\texttt{ring-prod-n-is-ring}]
  $(\forall i \in \mathit{NN}.\;\mathit{IS\text{-}RING}(f(i)))
  \;\wedge\; n \in \mathit{NN}
  \;\Rightarrow\; \mathit{IS\text{-}RING}(\mathit{RING\text{-}PROD\text{-}N}(f,n))$.\\
  (Derivable by $\mathit{NN}$-induction using \texttt{zero-ring-is-ring}
  and \texttt{ring-prod-is-ring}; installed as an axiom for direct use.)
\end{description}

\begin{remark}
The indexed-family product $\prod_{i \in I} f(i)$ for an arbitrary
index set $I$ (not necessarily $\mathit{NN}$) requires the big-union
constructor $\bigcup_{i \in I} \mathit{A}(f(i))$ to serve as the
function codomain.  This depends on \texttt{UNION-SET}, which is
currently pending.  The family case will be added once
\texttt{UNION-SET} is in place.
\end{remark}
